\documentclass[letterpaper,10pt]{article}
\usepackage[left=2cm,top=3cm,right=2cm,bottom=4cm]{geometry} 
\usepackage[utf8]{inputenc}
\usepackage[activeacute,spanish]{babel}
\usepackage{multirow}	% Para poder unir filas en las tablas
\usepackage{colortbl}	% Para colorear tablas
\usepackage{latexsym,amssymb,oldlfont,amsmath}
\usepackage{graphicx,graphics}
\usepackage{multicol}
\usepackage{color}
\usepackage{listings}
\usepackage{lastpage}
\usepackage{enumitem}



%\usepackage{fontspec}
%\setmainfont{Times New Roman}
\newcommand{\R}{{\mathbb R}}
\newcommand{\si}{{\mathcal S}^n}
\renewcommand{\theequation}{\thesection.\arabic{equation}} % Cambia formato salida numeracion equation
\makeatletter\@addtoreset {equation}{section} \makeatother % Tambien se usa
\renewcommand{\baselinestretch}{1.3}
\usepackage{amssymb}
\usepackage{hyperref}
\definecolor{midnightblue}{rgb}{0.1,0.1,0.44}


\newcommand{\dis}{\displaystyle}
\renewcommand{\theenumii}{\arabic{enumii}}

\def\Z{\mathbb{Z}}
\def\N{\mathbb{N}}
\def\R{\mathbb{R}}
\def\M{\mathcal{M}_n(\mathbb{R})}
\def\sin{\operatorname{sen}}
\def\e{\operatorname{e}}
\DeclareMathOperator{\arcsec}{arcsec}
\DeclareMathOperator{\arccot}{arccot}
\DeclareMathOperator{\arccsc}{arccsc}
\providecommand{\abs}[1]{\lvert#1\rvert}


\usepackage{etoolbox,fancyhdr,xcolor}
\fancyhf{} % sets both header and footer to nothing
\renewcommand{\headrulewidth}{0pt}
%\renewcommand{\footrulewidth}{0.5pt}
\pagestyle{fancy}
\headheight=60pt %para cambiar el tamaño del encabezadox
\fancyhead[L] %la "L" indica a la izquierda
{
\begin{minipage}{6cm}
\includegraphics[width=0.9\textwidth]{logo_udp.eps}
\end{minipage}
}

%\fancyhead[C]
%{
%\textbf{Cálculo II - Guía 1}
%}


%\fancyhead[R] %la "L" indica a la izquierda
%{

%}
%\begin{minipage}{4cm}
%\small %tamaño de letra pequeño
%{
%\textsf
%{
%Universidad de Chile\\
%...}}
%\end{minipage}}



%\fancyfoot[C]{}
%{
%\begin{minipage}{10cm}
%\color{midnightblue}
%\begin{center}Casa Matriz: Monjitas 527 Piso 15 - Santiago de Chile\\
%Fono (02)24367255 \end{center}
%\end{minipage}
%}
\fancyfoot[C]{}
%Página \thepage\ de \pageref{LastPage}}% Páginas Pares numeradas a la derecha
%\fancyfoot[LO]{\thepage} % Páginas Impares numeradas a la izquierda

\title{C2_Ay5}

\begin{document}

\begin{center}
Cálculo II\\
Ayudantía Nº5\\
Primer Semestre 2017
\end{center}



\begin{enumerate}
\item Determine si las siguientes afirmaciones son verdaderas o falsas:
\begin{enumerate}[label=\emph{\alph*)}]
\item $\dis \int_0^1 x\, dx \leq \int_0^1 x^2\, dx$
\item Si $f$ y $g$ son funciones integrables en todo intervalo $[a,b] \subset R$ tales que
\begin{center}
$\dis \int_1^2 f(x)\, dx=-1  \hspace{1cm}  \int_1^5 f(x)\, dx=3 \hspace{1cm} \dis \int_1^5 g(x)\, dx=1 $
\end{center}
Entonces $\dis \int_2^5 f(x)\, dx=4$ \hspace{2mm} y \hspace{2mm} $\dis \int_1^5 g(x)\left(\int_1^5 f(u)du\right)  \, dx=3$
\end{enumerate}
\vspace{2mm}


\item Calcular: \\
\[ \lim_{x \to 0} \frac{\ln(1+x^2)+\dis \int_1^{\cos x} e^{-t^2}dt}{xe^{x^3}+x^2-\dis \int_0^{x}e^{t^3}dt}\]
\vspace{2mm}

\item Sea $F$ la función definida por $\dis F(x)= \int_{\frac{1}{2}}^{x}f(u)du$, donde $\dis f(u)= \int_0^u \frac{1}{1+t^2}dt$. Determine $F'(1)$.
\vspace{2mm}

\item Sea $\dis f(x)= \int_{0}^{g(x)} \frac{dt}{\sqrt{1+t^3}}$, donde $\dis g(x)= \int_{0}^{\cos(x)} \left(1+ \sin (t^2)\right)dt$. Calcule $\dis f'\left(\frac{\pi}{2}\right)$.
\vspace{4mm}


\item 
\begin{enumerate}[label=\emph{\alph*)}]
\item Sean $\dis \int_1^{11} f(x)\, dx = 15$ y $\dis \int_1^{5} f(x)\, dx = 8$. Calcule $\dis \int_2^{4} f(3x-1)\, dx$
\item Calcule $\dis \int_{-1}^{1} \sqrt{\abs{x}-x} \, dx$
\end{enumerate}




\end{enumerate}




\vspace{10mm}




\end{document}
